\documentclass[11pt]{article}
\usepackage[T1]{fontenc}
\usepackage{textcomp}
\usepackage[margin=0.75in]{geometry}
\usepackage{amsmath,amssymb,amsthm}
\usepackage{array,booktabs}
\usepackage{hyperref}
\setlength{\parindent}{0pt}
\newtheorem{theorem}{Theorem}
\theoremstyle{definition}
\newtheorem{definition}{Definition}
\title{Flow Proof Artifact: prop-01-equimultiple-of-equimultiple}
\author{Generated by \texttt{flow doc proof}}
\date{Flow Proof Book}
\begin{document}
\maketitle
If any number of magnitudes are equimultiples of as many others, each of each, whatever multiple one is of one, that multiple the sum is of the sum.\\[0.5em]
\textbf{Source.} euclid — Elements, Book V, Proposition 1\\[0.5em]
\bigskip
\noindent{\large \textbf{Derived fact 1.} \textit{If any number of magnitudes are equimultiples of as many others, each of each, whatever multiple one is of one, that multiple the sum is of the sum} \hfill the Euclidean plane \cdot Euclid Book V \cdot Proposition 1: equimultiples of equimultiples are equimultiple of the originals}\par
\smallskip\noindent\textit{Coordinate: the Euclidean plane \cdot Euclid Book V \cdot Proposition 1: equimultiples of equimultiples are equimultiple of the originals}\par
\smallskip\noindent\textit{Source: euclid — Elements, Book V, Proposition 1}\par
\medskip
\noindent\textbf{Goal.} If any number of magnitudes are equimultiples of as many others, each of each, whatever multiple one is of one, that multiple the sum is of the sum\par
\noindent$\displaystyle \text{equimultiple of equimultiple is equimultiple}$\par
\medskip
\noindent
\begin{tabular}{@{} >{\raggedright\arraybackslash}p{0.06\textwidth} >{\raggedright\arraybackslash}p{0.47\textwidth} >{\raggedleft\arraybackslash}p{0.41\textwidth} @{}}
\textbf{\#} & \textbf{Proof} & \textbf{Mathematics} \\
\midrule
\textbf{1.} & We prove that proposition 1: equimultiples of equimultiples are equimultiple of the originals for Euclid Book V the Euclidean plane. & \\[0.45em]
\textbf{2.} & Let m = multiple\_factor. & \\[0.45em]
\textbf{3.} & Let a = first\_magnitude. & \\[0.45em]
\textbf{4.} & Let b = second\_magnitude. & \\[0.45em]
\textbf{5.} & From step 2, step 3, and step 4, we can deduce that equimultiple(m, sum) equals sum of equimultiples. & $\displaystyle equimultiple(m, sum) = \text{sum of equimultiples}$ \\[0.45em]
\textbf{6.} & We can deduce that equimultiple of equimultiple is equimultiple. Hence proven. & $\displaystyle \text{equimultiple of equimultiple is equimultiple}$ \\[0.45em]
\bottomrule
\end{tabular}
\label{thm:Geometry-Euclid Book V-proposition_1_equimultiples_of_equimultiples_are_equimultiple_of_the_originals}
\par\medskip\hrule
\end{document}
